\documentclass[12pt]{article}
\usepackage[T1]{fontenc}
\usepackage[utf8]{inputenc}
\usepackage{lmodern}
\usepackage{textcomp}
\usepackage{enumitem}
\usepackage{geometry}
\geometry{margin=1in}
\begin{document}
\begin{center}
Answer Key - History - K-12 Level Assessment 6 \
\small (Answers are ordered exactly as in the question paper.)
\end{center}

\section*{Multiple Choice}
\begin{enumerate}[label=\arabic*.]
\item \textbf{Answer:} B
\\ \textit{Options:} A) did not use pictographs; B) appears to have developed suddenly; C) was the earliest form of writing in the world; D) all of the above
\\ \textit{Note:} The Egyptian system of hieroglyphics is considered to have developed suddenly because archaeological evidence indicates that this complex writing system emerged fully formed without a clear antecedent, suggesting a rapid evolution rather than a gradual progression.
\item \textbf{Answer:} B
\\ \textit{Options:} A) emphasized flakes over blades.; B) changed rapidly.; C) involved elaborate preparation of core tools.; D) all of the above
\\ \textit{Note:} The Upper Paleolithic period is characterized by significant advancements in tool technology, including the development of more specialized and efficient tools, which reflects a rapid evolution compared to the relatively static tool technologies of the Middle Paleolithic.
\item \textbf{Answer:} A
\\ \textit{Options:} A) When sled dogs were first domesticated; B) When the Bering Strait was exposed and open for travel; C) When eastern Siberia was first inhabited; D) The age of the earliest New World sites
\\ \textit{Note:} Anthropologists do not need to know when sled dogs were first domesticated because the entry of people into the Americas from Siberia is primarily determined by archaeological and genetic evidence rather than the domestication of specific animal species.
\item \textbf{Answer:} D
\\ \textit{Options:} A) is made out of stone, wood, or bone.; B) needs constant cleaning and maintenance.; C) requires metal and a source of heat to be forged; D) requires many more flakes to be removed.
\\ \textit{Note:} A hand axe is a more complex and advanced tool than an Oldowan tool, necessitating the removal of many more flakes to achieve its refined shape and sharp edges.
\item \textbf{Answer:} D
\\ \textit{Options:} A) trade.; B) warfare.; C) the development of a common religion.; D) conspicuous consumption by elites.
\\ \textit{Note:} The correct answer is based on the observation that Minoan culture, unlike many other early civilizations, did not exhibit clear signs of wealth disparity or extravagant displays of luxury among elites, suggesting a more egalitarian social structure.
\end{enumerate}

\section*{True / False}
\begin{enumerate}[label=\arabic*.]
\setcounter{enumi}{5}
\item \textbf{Answer:} True
\\ \textit{Options:} A) True; B) False
\\ \textit{Note:} The term "prognathous" specifically describes a condition where the lower jaw is positioned forward in relation to the upper jaw, affecting facial structure.
\item \textbf{Answer:} True
\\ \textit{Options:} A) True; B) False
\\ \textit{Note:} The statement is true because during the late Miocene, dramatic climate shifts led to habitat changes that resulted in the extinction of many ape species that could not adapt to the new environmental conditions.
\item \textbf{Answer:} True
\\ \textit{Options:} A) True; B) False
\\ \textit{Note:} The peak of Minoan civilization occurred around 1600 BCE, which is historically linked to the catastrophic eruption of Santorini and a significant earthquake that impacted their society.
\item \textbf{Answer:} True
\\ \textit{Options:} A) True; B) False
\\ \textit{Note:} Ancestral Pueblo elites gained power by controlling food surpluses, which allowed them to provide for their communities during times of scarcity and thus establish their authority and influence.
\item \textbf{Answer:} True
\\ \textit{Options:} A) True; B) False
\\ \textit{Note:} Mastabas, which are rectangular, flat-roofed structures made of mud-brick, served as the earliest form of tombs for the elite in ancient Egypt and laid the architectural foundation that eventually developed into the more complex pyramid structures.
\end{enumerate}

\section*{Long Form Response}
\begin{enumerate}[label=\arabic*.]
\setcounter{enumi}{10}
\item \textbf{Answer:} The Hassunan and Samarran sites in Mesopotamia both relied heavily on agriculture, which played a crucial role in shaping their societies. These early agricultural practices included the cultivation of staple crops such as wheat and barley, supported by advanced irrigation techniques that utilized the region's rivers. This reliance on farming allowed both cultures to develop surplus food, which in turn facilitated population growth and the establishment of more complex social structures. Additionally, the successful agricultural output contributed to trade and interactions with neighboring communities, further enriching their societies. Overall, their agricultural practices were foundational to their economic stability and cultural development.
\\ \textit{Note:} correct because it accurately identifies the shared agricultural practices of the Hassunan and Samarran sites, emphasizing the cultivation of staple crops and the use of irrigation, which together fostered food surpluses and contributed to societal development.
\item \textbf{Answer:} The Olmec civilization, often considered the "mother culture" of Mesoamerica, held complex religious beliefs centered around nature, spirituality, and the idea of a cosmic balance. They worshipped a pantheon of deities, often depicted in their art and iconography, which included colossal stone heads and intricate jade carvings that symbolized important figures and natural elements. These artistic representations not only served religious purposes but also established a shared cultural identity across different Olmec territories, reinforcing their beliefs and values. The consistent themes in their art, such as the jaguar and the concept of duality, helped unify diverse communities by promoting a common understanding of their spiritual world. Through their art and symbols, the Olmec effectively communicated their religious beliefs, fostering connections among people across vast regions.
\\ \textit{Note:} correct because it accurately identifies the Olmec civilization's religious beliefs related to nature and spirituality, highlights the significance of their art and iconography in expressing these beliefs, and emphasizes how these cultural elements fostered a shared identity across different regions.
\end{enumerate}

\vfill
\noindent\textit{End of Answer Key}
\end{document}